\documentclass[11pt,a4paper]{article}

\usepackage[T1]{fontenc}
\usepackage{lmodern}
\usepackage[utf8]{inputenc}
\usepackage{textcomp}

\usepackage{amsmath,amssymb,amsthm}
\usepackage{mathtools}
\usepackage{hyperref}
\usepackage{doi}
\usepackage{geometry}
\usepackage{booktabs}
\usepackage{array}
\usepackage{enumitem}
\usepackage{xcolor}

\geometry{margin=2.5cm}

\input{glyphtounicode}
\pdfgentounicode=1

\hypersetup{
  colorlinks=true,
  linkcolor=black,
  citecolor=black,
  urlcolor=black,
  pdftitle={Inverse-Square Critical Coverage and Constant Saturation Depth
    from an Exact Interval Theorem},
  pdfauthor={J\'{e}r\^{o}me Beau},
  pdfsubject={Exact asymptotics of Weil--BFS critical coverage on finite Heisenberg graphs},
  pdfkeywords={critical coverage, BFS saturation depth, Heisenberg group, Weil representation,
    Fourier character, interval theorem, spectral admissibility}
}

\theoremstyle{plain}
\newtheorem{theorem}{Theorem}
\newtheorem{lemma}{Lemma}
\newtheorem{corollary}{Corollary}

\theoremstyle{definition}
\newtheorem{definition}{Definition}
\newtheorem{proposition}{Proposition}
\newtheorem{remark}{Remark}

\newcommand{\Heis}{\mathrm{Heis}}
\newcommand{\Cay}{\mathrm{Cay}}
\newcommand{\BFS}{\mathrm{BFS}}
\newcommand{\Bball}{\lvert B_n\rvert}
\newcommand{\Cheis}{C_{\mathrm{Heis}}}
\newcommand{\epssat}{\varepsilon_{\mathrm{sat}}}
\newcommand{\xone}{x_1}
\newcommand{\none}{n_1}
\newcommand{\Tp}{T'_c}
\newcommand{\Fq}{\mathbb{Z}/q\mathbb{Z}}

\title{Inverse-Square Critical Coverage and Constant Saturation Depth\\
from an Exact Interval Theorem\\[0.4em]
\large Weil--BFS Novelty on Heisenberg Graphs\\[0.2em]
\large Cosmochrony --- Spectral Admissibility Sub-programme}

\author{J\'{e}r\^{o}me Beau\\[0.3em]
\small Independent Researcher\\
\small \texttt{jerome.beau@cosmochrony.org}}

\date{2026}

\begin{document}

  \maketitle

  \begin{abstract}
    The auto-calibrated saturation depth $\none(q)$ for BFS on $\Heis_3(\Fq)$ --- the central
    asymptotic variable of the O25--O28 window
    calibration --- converges in probability to the constant $n_* = 22$ under uniform generic
    block sampling, and the critical coverage
    $\xone(q) = \lvert B_{\none}\rvert/q^2$ obeys the inverse-square law $\xone(q) \asymp q^{-2}$.
    The mechanism is an exact algebraic reduction.
    Every implemented three-step Weil fingerprint is a non-zero scalar multiple of a single Fourier
    character, so the cumulative rank is a sumset cardinality,
    \[
      r_c(n) = \bigl\lvert \Tp + C\,I_{n+1} \bigr\rvert,
      \qquad
      \Tp = \{(c_2+c_3)\eta_2 + c_3\eta_3 : \eta_i \in \{-1,0,1\}\},
    \]
    with $C = c_1+c_2+c_3$ and $I_m$ the reduction of $[-m,m]$; the shellwise novelty obeys the
    universal deterministic ceiling $\Delta r_c(n) \le 2\lvert\Tp\rvert \le 18$.
    Exact sphere counts of the infinite discrete Heisenberg group give the strict threshold crossing
    $s_{22} = 16{,}934 < 18/\epssat = 18{,}000 < s_{23} = 19{,}412$.
    Deterministically, for every prime $q \ge 311$ and every admissible block sequence,
    $\none(q) \le 22$ and $5/q^2 \le \xone(q) \le 99{,}689/q^2$; under uniform generic block sampling,
    $\none(q) \to 22$ and $q^2\,\xone(q) \to \lvert B_{22}\rvert = 99{,}689$ in probability.
    The measured depths ($\none = 14$ at $q=211$, $19$ at $q=601$) sit in a transitional regime
    governed by block resonances, whose deficits produce the observed effective exponent
    $\approx -0.76$ en route to the exact asymptotic $-2$.
    A structural corollary is that the implemented auto-calibrated fitting window closes
    deterministically for $q > 937{,}024$.
    As interpretation (structural reading, not a result): the finite Weil--BFS observable has a
    finite novelty horizon --- beyond a fixed relational depth, independent of the arithmetic size
    of the substrate, exploration creates no above-threshold spectral novelty per vertex --- so
    shellwise novelty is a local combinatorial phenomenon, and no continuum-scaling channel opens
    from this observable.
  \end{abstract}

  \medskip
  \noindent\textbf{Keywords.}
  critical coverage; BFS saturation depth; interval theorem; sumset growth; marginal novelty; Weil
  representation; Heisenberg group; Fourier character; projective capacity; spectral admissibility.

  \newpage
  \tableofcontents

  \newpage

  \section{The saturation-depth problem}
    \label{sec:problem}

    Let $q$ be an odd prime and $G_q = \Heis_3(\Fq)$, with multiplication
    \[
      (a,b,\gamma)(a',b',\gamma') = (a+a',\,b+b',\,\gamma+\gamma'+ab').
    \]
    Equip $G_q$ with the Cayley graph $\Cay(G_q,S_q)$ for
    $S_q = \{X^{\pm 1},Y^{\pm 1}\}$, where $X = (1,0,0)$ and $Y = (0,1,0)$.
    Breadth-first search from the identity produces spheres $S_n$ and balls
    $B_n = \bigcup_{m \le n} S_m$.
    For a block multi-index $c = (c_1, c_2, c_3) \in (\Fq)^3$, each shell element $g \in S_n$ and
    each path $w = (s_1, s_2, s_3) \in S_q^3$ yield a $k=3$ Weil fingerprint vector in
    $\mathbb{C}^q$: the pointwise product of the three vectors $\rho_{c_j}(g s_1 \cdots s_j)\,u$,
    where $u$ is the uniform seed and $\rho_c$ is the $q$-dimensional Schr\"{o}dinger
    representation of the finite Heisenberg group used in the O12 Weil-block construction,
    $(\rho_c(a,b,\gamma)f)(x) = \psi_c(\gamma + bx)\,f(x+a)$, $\psi_c(t) = e^{2\pi i c t/q}$, of
    O12~\cite{Beau2026a16}.
    The cumulative Gram--Schmidt rank $r_c(n) = \sum_{m \le n}\Delta r_c(m)$ is monotone and bounded
    by $q$, and the capacity-novelty density is
    \begin{equation}
      \sigma_c(n) = \frac{\Delta r_c(n)}{\lvert S_n\rvert}.
      \label{eq:sigma}
    \end{equation}
    The auto-calibrated fitting window of O12~\cite{Beau2026a16} has upper edge
    \begin{equation}
      \none(q) = \text{last } n \le N(q) \text{ with } \bar{\sigma}(n) > \epssat,
      \qquad \epssat = 10^{-3},
      \label{eq:n1def}
    \end{equation}
    where $\bar\sigma$ is the mean of~\eqref{eq:sigma} over $M = 5$ probe blocks and
    $N(q) = \lceil\sqrt{q}\rceil + 2$ is the processing cap of the production measurement.

    \begin{remark}[The depth is a marginal-novelty stop, not a filling time]
      \label{rem:marginal}
      At the window edge the cumulative rank is far from full: $r(\none) \approx 137/307$ at $q=307$
      and $183/601$ at $q=601$.
      So $\none$ is \emph{not} the depth at which the fingerprints span $\mathbb{C}^q$.
      By~\eqref{eq:n1def} it is the depth at which the marginal rank gained per explored vertex,
      $\sigma(n)$, falls below the fixed threshold $\epssat$: a stop of marginal novelty.
    \end{remark}

    O25--O28~\cite{Beau2026a29,Beau2026a32} identified the asymptotic law of $\none(q)$ as the
    central open direction of the window calibration.
    O28 rejected linear scaling out of sample and found a decreasing ratio $\none(q)/q$ over the
    measured primes.
    This note determines the law completely: $\none(q)$ is asymptotically \emph{constant}, and the
    natural coverage observable decays as an exact inverse square.

    Two background facts about the graph are used throughout.

    \begin{proposition}[Volume growth, $D=4$]
      \label{prop:growth}
      The infinite discrete Heisenberg group with the standard generators has homogeneous dimension
      $D = 4$ (Bass--Guivarc'h~\cite{Bass1972,Guivarch1973}; O9~\cite{Beau2026a13}), so
      \begin{equation}
        \lvert B_n\rvert = \Cheis\,n^4\,(1 + o(1))
        \qquad (n\to\infty),
        \qquad \Cheis \approx 0.427,
        \label{eq:growth}
      \end{equation}
      with shell ratio $\lvert S_n\rvert / n^3 \to 4\Cheis \approx 1.71$.
      Finite-group balls agree with these infinite-group balls in the injective range of
      Lemma~\ref{lem:isometry} below.
    \end{proposition}

    \begin{remark}[Diffusive spectral gap (background only)]
      \label{prop:gap}
      A separate numerical diagnostic finds
      $\lambda_2(L_{G_q}) \approx 4\pi^2/q^2$, the gap of the horizontal abelianisation torus
      $(\mathbb{Z}/q\mathbb{Z})^2$.
      Theorem~\ref{thm:interval} below shows that this relaxation scale plays no role in the rank
      dynamics: the fingerprint rank is a sumset cardinality, blind to mixing.
    \end{remark}

    \begin{definition}[Critical coverage]
      \label{def:x1}
      \begin{equation}
        \xone(q) = \frac{\lvert B_{\none(q)}\rvert}{q^2}.
        \label{eq:x1}
      \end{equation}
    \end{definition}

    \begin{proposition}[Exact depth--coverage identity]
      \label{prop:reduction}
      For $n\ge1$, set $\kappa_n=\lvert B_n\rvert/n^4$.
      Then, exactly,
      \begin{equation}
        \none(q) = \left(\frac{\xone(q)}{\kappa_{\none(q)}}\right)^{1/4}\sqrt{q}.
        \label{eq:reduction}
      \end{equation}
      In the generic limit below, $\kappa_{\none(q)}\to\kappa_{22}=99{,}689/22^4$ in probability.
      The asymptotic constant $\Cheis$ is background geometry, not a substitute for this exact
      finite-depth coefficient.
    \end{proposition}

  \section{The Fourier-character reduction}
    \label{sec:fourier}

    \begin{proposition}[Fingerprints are characters]
      \label{prop:character}
      Let $g \in G_q$, $w = (s_1,s_2,s_3) \in S_q^3$, and $g_j = g s_1 \cdots s_j =
      (a_j, b_j, \gamma_j)$.
      The $k=3$ fingerprint of $(g,w)$ in block $c$ is
      \begin{equation}
        f_{g,w,c}(x)
        = q^{-3/2}\,A(g,w,c)\,
          \exp\!\Bigl(\tfrac{2\pi i}{q}\Bigl(\textstyle\sum_{j=1}^{3}c_j b_j\Bigr)x\Bigr),
        \qquad \lvert A(g,w,c)\rvert = 1.
        \label{eq:character}
      \end{equation}
      The coordinates $a_j$ and the central coordinates $\gamma_j$ enter only the scalar
      $A(g,w,c)$.
      Consequently distinct frequencies are exactly orthogonal, and the cumulative Gram--Schmidt
      rank equals the cardinality of the cumulative frequency set: no numerical rank tolerance is
      mathematically relevant.
    \end{proposition}

    \begin{proof}
      Each factor is $(\rho_{c_j}(a_j,b_j,\gamma_j)u)(x) = q^{-1/2}\psi_{c_j}(\gamma_j + b_j x)$,
      since the uniform seed is translation invariant.
      The pointwise product of the three factors collects the constant phases
      $\psi_{c_j}(\gamma_j)$ into $A$ and the linear phases into the single frequency
      $\sum_j c_j b_j$.
    \end{proof}

    \begin{remark}[Canonical block projection]
      \label{rem:projection}
      The pointwise product realises the block projection of O12 exactly and canonically: the
      Hadamard map $m(f_1 \otimes f_2 \otimes f_3)(x) = f_1(x)f_2(x)f_3(x)$ intertwines the
      diagonal action $\rho_{c_1} \otimes \rho_{c_2} \otimes \rho_{c_3}$ with $\rho_{C}$,
      $C = c_1+c_2+c_3 \ne 0$, and is a coisometry onto the diagonal $q$-dimensional constituent
      $\mathrm{span}\{e_x^{\otimes 3}\}$ of the triple tensor product; see
      Gurevich--Hadani~\cite{GurevichHadani2007} for the geometric background on finite Weil
      representations.
      A structural consequence is that the rank observable is blind to the central coordinate: with
      the uniform seed, $\gamma$ and $a$ contribute unimodular scalars only.
      Sensitivity to the central extension would require a different seed or a matrix-coefficient
      observable, which would define a different, separately calibrated problem.
    \end{remark}

  \section{The exact interval theorem}
    \label{sec:interval}

    Write $\pi_q:\mathbb{Z}\to\Fq$ for the reduction and $I_m = \pi_q([-m,m])$.
    For a block $c$ with $C = c_1+c_2+c_3 \ne 0$ define
    \begin{equation}
      \Tp = \{(c_2+c_3)\eta_2 + c_3\eta_3 : \eta_2,\eta_3 \in \{-1,0,1\}\}
      \subseteq \Fq,
      \qquad 0 \in \Tp = -\Tp,
      \quad \lvert\Tp\rvert \le 9.
      \label{eq:Tprime}
    \end{equation}

    \begin{lemma}[$b$-marginal of the ball]
      \label{lem:bmarg}
      For every $n \ge 0$, the set of $b$-coordinates of $B_n$ is exactly $I_n$; in particular it
      has $\min(2n+1, q)$ elements.
    \end{lemma}

    \begin{proof}
      The map $\beta:(a,b,\gamma)\mapsto b$ is a homomorphism to $(\Fq,+)$.
      Moreover,
      \[
        \beta(S_q)=\{\pi_q(-1),0,\pi_q(1)\}.
      \]
      Thus $\beta(B_n) \subseteq I_n$; conversely $Y^k \in B_n$ for
      $\lvert k\rvert \le n$.
    \end{proof}

    \begin{theorem}[Exact interval formula]
      \label{thm:interval}
      For every prime $q$, every block with $C \ne 0$, and every $n \ge 0$,
      \begin{equation}
        r_c(n) = \bigl\lvert\, \Tp + C\,I_{n+1} \,\bigr\rvert.
        \label{eq:interval}
      \end{equation}
    \end{theorem}

    \begin{proof}
      By Proposition~\ref{prop:character} the rank is the cardinality of the cumulative frequency
      set $F(n)$.
      With $\eta_i = \beta(s_i) \in \{-1,0,1\}$ and $b_j = b + \eta_1 + \cdots + \eta_j$,
      \begin{equation}
        \textstyle\sum_{j=1}^{3} c_j b_j
        = C\,b + C\eta_1 + (c_2+c_3)\eta_2 + c_3\eta_3 .
        \label{eq:frequency}
      \end{equation}
      All $27$ values of $(\eta_1,\eta_2,\eta_3)$ are realised at every start element (each
      $\eta$-value is realised by a generator), and by Lemma~\ref{lem:bmarg} every $b \in I_n$ is
      realised in $B_n$.
      Hence $F(n) = C\,I_n + C\{-1,0,1\} + \Tp = C\,I_{n+1} + \Tp$, using
      $I_n + I_1 = I_{n+1}$.
    \end{proof}

    The coefficient of $\eta_1$ in~\eqref{eq:frequency} is exactly $C$, the step of the frequency
    progression: the first path increment is absorbed into the same progression as the ball
    coordinate.
    The $27$ nominal centres therefore act through at most $9$ progression classes, which is the
    content of the following universal bound.

    \begin{corollary}[Universal deterministic bounds]
      \label{cor:bounds}
      For every prime $q$ and every block with $C \ne 0$:
      \begin{enumerate}[nosep]
        \item $\Delta r_c(n) \le 2\lvert\Tp\rvert \le 18$ for every $n \ge 1$, and
          $r_c(0) \le 27$;
        \item $\Delta r_c(n) \ge 2$ whenever $r_c(n-1) \le q-2$.
      \end{enumerate}
    \end{corollary}

    \begin{proof}
      (i) $F(n)\setminus F(n-1) \subseteq (\Tp + C\pi_q(n+1)) \cup (\Tp - C\pi_q(n+1))$.
      (ii) Multiplying by $C^{-1}$, $C^{-1}F(n) = W + \{-1,0,1\}$ with $W = C^{-1}\Tp + I_n$ a
      union of maximal cyclic arcs; adding $\{-1,0,1\}$ extends every maximal arc by one point on
      each side, so the growth is at least $2$ unless the complement of $W$ is a single point.
    \end{proof}

    Part (ii) explains the slow blocks observed in the O25 pipeline~\cite{Beau2026a29}, which add
    exactly two directions per shell: they are the fully resonant blocks of the classification
    below.

    \begin{theorem}[Shellwise law and resonances]
      \label{thm:shellwise}
      Let $q\ge11$, fix a block with $C \ne 0$, and let $1 \le n \le (q-3)/2$.
      Define the blocked set
      $R(n) = \{t \in \Tp : \exists\,t' \in \Tp,\ \exists\,m \in [1, 2n+1],\ t'-t = C\pi_q(m)\}$
      and the transient count
      $\tau_c(n) = \#\{\{t,t''\} \subseteq \Tp : t''-t = 2C\pi_q(n+1),\ t \notin R(n),\
      t'' \notin -R(n)\}$.
      Then
      \begin{equation}
        \Delta r_c(n) = 2\bigl(\lvert\Tp\rvert - \lvert R(n)\rvert\bigr) - \tau_c(n).
        \label{eq:shellwise}
      \end{equation}
      Moreover:
      \begin{enumerate}[nosep]
        \item a stable resonance $t'-t = C\,m$ activates at every depth
          $n \ge \lceil (m-1)/2 \rceil$ and its deficit is even and non-decreasing in $n$;
        \item $\lvert\Tp\rvert < 9$ iff $(c_2+c_3)/c_3 \in \{0, \pm 1, \pm 2, \pm\tfrac12\}$ in
          $\mathbb{F}_q$, with $\lvert\Tp\rvert \in \{3, 5, 7\}$ in the respective degenerate
          cases;
        \item each degeneracy, stable-resonance, or transient condition is a non-trivial linear
          condition on $(c_1,c_2,c_3)$, so under uniform generic sampling each has probability at
          most $1/(q-4)$.
      \end{enumerate}
    \end{theorem}

    \begin{proof}
      The right-edge candidate $t + C\pi_q(n+1)$ lies in $F(n-1) = \Tp + C I_n$ iff
      $t'-t = C\pi_q(n+1-j)$ for some $j \in [-n,n]$, i.e.\ iff $t \in R(n)$ with
      $m = n+1-j \in [1, 2n+1]$; the left edge is the mirror case under $\Tp = -\Tp$, so the two
      blocked counts are equal.
      A right candidate coincides with a left candidate iff $t''-t = 2C\pi_q(n+1)$, and each such
      pair with both members new is counted once, giving $-\tau_c(n)$.
      For (ii), a collision $(c_2+c_3)\Delta_2 + c_3\Delta_3 \equiv 0$ with
      $(\Delta_2,\Delta_3) \ne (0,0)$, $\lvert\Delta_i\rvert \le 2$, forces the stated ratio; each
      degenerate ratio yields the stated cardinality by direct enumeration.
      For (iii), each condition is a non-trivial linear form on $(c_1,c_2,c_3)$ (for stable and
      transient resonances the coefficient of $c_1$ is $m \ne 0$), hence has at most $q^2$
      solutions, while the generic set has $(q-1)(q^2-3q+3) \ge q^2(q-4)$ elements.
    \end{proof}

  \section{Exact spheres and the limiting depth}
    \label{sec:spheres}

    The limiting depth is governed by the sphere sizes $s_n$ of the \emph{infinite} discrete
    Heisenberg group $H_3(\mathbb{Z})$ with the same generators, whose exact word metric is given
    by Blach\`{e}re~\cite{Blachere2003}.
    The values below are exact integer BFS counts of $H_3(\mathbb{Z})$ to depth $31$.

    \begin{table}[ht]
      \centering
      \caption{Exact sphere and ball sizes of $H_3(\mathbb{Z})$ near the threshold, and the
        novelty bound $18/s_n$ against $\epssat = 10^{-3}$.
        The crossing between $n = 22$ and $n = 23$ is strict, with margins $+6.3\%$ and $-7.3\%$.}
      \label{tab:spheres}
      \begin{tabular}{rrrr}
        \toprule
        $n$ & $s_n$ & $\lvert B_n\rvert$ & $18/s_n$ \\
        \midrule
        19 & 10{,}780 & 55{,}453  & $1.670 \times 10^{-3}$ \\
        20 & 12{,}626 & 68{,}079  & $1.426 \times 10^{-3}$ \\
        21 & 14{,}676 & 82{,}755  & $1.226 \times 10^{-3}$ \\
        22 & 16{,}934 & 99{,}689  & $1.063 \times 10^{-3}$ \\
        23 & 19{,}412 & 119{,}101 & $0.927 \times 10^{-3}$ \\
        24 & 22{,}124 & 141{,}225 & $0.814 \times 10^{-3}$ \\
        \bottomrule
      \end{tabular}
    \end{table}

    Early values are $s_1,\dots,s_5 = 4, 12, 36, 82, 164$, matching the production BFS; beyond the
    table, $s_{25},\dots,s_{31} = 25{,}076$, $28{,}280$, $31{,}748$, $35{,}486$, $39{,}508$,
    $43{,}826$, $48{,}444$.
    On $[1,31]$ the sequence $s_n$ is strictly increasing and never equals $18{,}000$, so
    \begin{equation}
      n_*(\epssat) = \max\{n : 18/s_n > \epssat\} = 22,
      \qquad
      K = \lvert B_{22}\rvert = 99{,}689 .
      \label{eq:nstar}
    \end{equation}

    \begin{lemma}[Central extent and ball isometry]
      \label{lem:isometry}
      In $H_3(\mathbb{Z})$, $\max_{g \in B_n}\lvert\gamma(g)\rvert = \lfloor n^2/4\rfloor$.
      If $q > \max(2n,\, 2\lfloor n^2/4\rfloor)$, the reduction $H_3(\mathbb{Z}) \to G_q$ is
      injective and distance-preserving on $B_n$, so $s_m(q) = s_m$ and
      $\lvert B_m(q)\rvert = \lvert B_m\rvert$ for all $m \le n$.
    \end{lemma}

    \begin{proof}
      $\gamma$ changes only on $Y^{\pm}$-steps, by the current value of $a$; with $p$ $X$-steps
      and $r$ $Y$-steps, $p + r \le n$, one gets $\lvert\gamma\rvert \le pr \le \lfloor
      n^2/4\rfloor$, attained by $X^{\lfloor n/2\rfloor}Y^{n - \lfloor n/2\rfloor}$.
      If $g \ne g' \in B_n$ reduce to the same element, $q$ divides $a-a'$ with
      $\lvert a-a'\rvert \le 2n < q$, so $a = a'$, likewise $b = b'$, and then $q$ divides
      $\gamma-\gamma'$ with $\lvert\gamma-\gamma'\rvert \le 2\lfloor n^2/4\rfloor < q$: a
      contradiction.
      If some word of length $\ell < d_{\mathbb{Z}}(g)$ evaluated to the reduction of $g$ in
      $G_q$, its evaluation in $H_3(\mathbb{Z})$ would lie in $B_\ell \subseteq B_n$ and reduce to
      the same element while being distinct from $g$, contradicting injectivity.
    \end{proof}

    \begin{lemma}[Cubic sphere lower bound]
      \label{lem:cubic}
      In $H_3(\mathbb{Z})$, $s_n \ge \tfrac{2}{3}(n^3 - n)$ for every $n \ge 2$.
      Consequently $s_n > 18{,}000$ for every $n \ge 23$ (exact BFS for $23 \le n \le 31$, the
      bound for $n \ge 31$).
    \end{lemma}

    \begin{proof}
      Fix $a, b \ge 1$ with $a + b = n$ and signs $\epsilon, \delta \in \{\pm 1\}$.
      A word in $X^{\epsilon}, Y^{\delta}$ with $a$ $X$-steps and $b$ $Y$-steps evaluates to
      $(\epsilon a, \delta b, \epsilon\delta k)$ where $k = \sum_{i=1}^{b} x_i$ and
      $0 \le x_1 \le \cdots \le x_b \le a$ counts the $X$-steps preceding the $i$-th $Y$-step.
      Every $k \in [0, ab]$ is attained (partitions in an $a \times b$ box), and the
      abelianisation bound gives distance exactly $n$, so
      $s_n \ge 4\sum_{a=1}^{n-1}\bigl(a(n-a)+1\bigr) \ge \tfrac{2}{3}(n^3-n)$; at $n = 31$ this is
      $19{,}840 > 18{,}000$, and the bound is increasing.
    \end{proof}

  \section{Main theorem}
    \label{sec:main}

    \begin{theorem}[Constant saturation depth and inverse-square coverage]
      \label{thm:main}
      Let $q\ge311$ be prime, $\epssat = 10^{-3}$, $M = 5$, and let $\none(q)$ be the
      statistic~\eqref{eq:n1def} with cap $N(q) = \lceil\sqrt{q}\rceil + 2$.
      \begin{enumerate}[nosep]
        \item \textbf{Deterministic part.}
          For arbitrary blocks with $C_i \ne 0$:
          \begin{equation}
            \none(q) \le 22,
            \qquad
            \frac{5}{q^2} \le \xone(q) \le \frac{99{,}689}{q^2},
          \end{equation}
          so $\xone(q) \asymp q^{-2}$ and $\xone(q) \to 0$ deterministically; in particular the
          positive-edge scenario ($\xone \ge c > 0$) and every logarithmic law are excluded.
        \item \textbf{Probabilistic part.}
          If the $M$ blocks are drawn independently and uniformly from the generic set
          $\{c \in (\mathbb{F}_q^{\times})^3 : C \ne 0\}$, then
          \begin{equation}
            \mathbb{P}\bigl(\none(q) = 22\bigr) \ge 1 - \frac{5640}{q-4},
          \end{equation}
          hence $\none(q) \to 22$ and $q^2\,\xone(q) \to K = 99{,}689$ in probability.
      \end{enumerate}
    \end{theorem}

    \begin{proof}
      (i) For $22 < n \le N(q)$, Corollary~\ref{cor:bounds} gives
      $\bar\sigma(n) \le 18/s_n(q)$.
      The cap satisfies $2N(q) \le 2\sqrt q+6<q$ and
      $2\lfloor N(q)^2/4\rfloor \le q/2+3\sqrt q+9/2<q$ for $q \ge 311$, so
      Lemma~\ref{lem:isometry} identifies $s_n(q) = s_n$ on the whole computed range, and
      $s_n > 18{,}000$ for every $n \ge 23$ by Lemma~\ref{lem:cubic}; hence
      $\bar\sigma(n) < \epssat$ throughout $(22, N(q)]$ and $\none \le 22$.
      Since $r_c(0) \le 27 \le q-2$, Corollary~\ref{cor:bounds} gives $\Delta r_c(1) \ge 2$ and
      $\bar\sigma(1) \ge 1/2 > \epssat$, so $\none \ge 1$ and
      $\lvert B_{\none}(q)\rvert \ge \lvert B_1(q)\rvert = 5$; on the other side,
      $\lvert B_{\none}(q)\rvert \le \lvert B_{22}(q)\rvert = \lvert B_{22}\rvert = 99{,}689$ by
      ball monotonicity and Lemma~\ref{lem:isometry} (which applies at depth $22$ since
      $q \ge 311 > 2\lfloor 22^2/4 \rfloor = 242$).
      (ii) Call a block clean if $\lvert\Tp\rvert = 9$, $R(22) = \emptyset$, and
      $\tau_c(22) = 0$.
      If all $M$ blocks are clean, Theorem~\ref{thm:shellwise} gives $\Delta r_{c_i}(22) = 18$ and
      $\bar\sigma(22) = 18/16{,}934 > \epssat$, so $\none = 22$ by part (i).
      Cleanliness fails on at most $24 \times 47 = 1128$ linear conditions per block (stable
      resonances with $m \le 45$, the transient $m = 46$, and the degeneracies), each of
      probability at most $1/(q-4)$ by Theorem~\ref{thm:shellwise}(iii); a union bound over $M=5$
      blocks gives the constant $5640$.
    \end{proof}

    \begin{remark}[The probability bound is conservative]
      \label{rem:conservative}
      The bound $1 - 5640/(q-4)$ is asymptotic and deliberately coarse: it is non-informative for
      $q \le 5644$ and weak well beyond, which does not affect the convergence since
      $5640/(q-4) \to 0$.
      Moreover $\none = 22$ only requires $\sum_i \Delta r_{c_i}(22) > \epssat M s_{22} = 84.67$,
      i.e.\ a total five-block deficit at depth $22$ of at most $5$, so the clean event is
      sufficient but far from necessary.
    \end{remark}

    \begin{remark}[Threshold stability]
      \label{rem:threshold}
      For every fixed $\varepsilon \in (0,1)$ the same proof gives
      $\none(q) \to n_*(\varepsilon) = \max\{n : 18/s_n > \varepsilon\}$ and
      $q^2 \xone(q) \to \lvert B_{n_*(\varepsilon)}\rvert$: the qualitative class --- constant
      depth, inverse-square coverage --- is threshold-independent; only the derived constant
      changes.
    \end{remark}

    \begin{corollary}[The auto-calibrated window closes at large $q$]
      \label{cor:window}
      The implemented lower window edge of O12~\cite{Beau2026a16} is the proxy
      $n_0(q)=\max\{2,\lceil(q/4)^{1/4}\rceil\}$.
      For $q > 4\cdot22^4 = 937{,}024$, $n_0(q) \ge 23 > 22 \ge \none(q)$: the auto-calibrated
      fitting window $[n_0, \none]$ is deterministically empty.
      Window-based exponent extraction on this observable is therefore confined to
      $q \le 937{,}024$.
    \end{corollary}

  \section{The transitional regime and the measured coverage}
    \label{sec:data}

    Table~\ref{tab:x1} reports the measured $\none(q)$ (auto-calibrated as in~\eqref{eq:n1def},
    reproducing the O28 production values exactly on $q \le 211$) and the exact coverage
    $\xone(q)$ via~\eqref{eq:x1}.

    \begin{table}[ht]
      \centering
      \caption{Measured saturation depth $\none(q)$, the ratio $\none/q$, and the critical
        coverage $\xone(q) = \lvert B_{\none(q)}\rvert/q^2$ using exact infinite-Heisenberg ball
        counts (the finite balls are isometric at every listed depth).
        Values $q \le 211$ reproduce O28~\cite{Beau2026a32}.}
      \label{tab:x1}
      \begin{tabular}{rrrr}
        \toprule
        $q$ & $\none(q)$ & $\none(q)/q$ & $\xone(q)$ \\
        \midrule
        29  & 4  & 0.138 & 0.161 \\
        61  & 8  & 0.131 & 0.482 \\
        101 & 11 & 0.109 & 0.616 \\
        151 & 13 & 0.086 & 0.535 \\
        211 & 14 & 0.066 & 0.368 \\
        307 & 16 & 0.052 & 0.296 \\
        401 & 17 & 0.042 & 0.221 \\
        503 & 19 & 0.038 & 0.219 \\
        601 & 19 & 0.032 & 0.154 \\
        \bottomrule
      \end{tabular}
    \end{table}

    The measured depths lie below the limit $22$ and climb toward it, and
    Theorem~\ref{thm:shellwise} identifies the mechanism: at accessible $q$ the expected number of
    active stable-resonance conditions per block at depth $n$ is of order $24(2n+1)/q = O(1)$, so
    resonance deficits keep $\Delta r$ below its ceiling $18$ and trigger the crossing early.
    As $q$ grows the resonance conditions thin out at rate $O(n/q)$, the deficits vanish, and
    $\none$ locks onto $22$; the coarse bound of Theorem~\ref{thm:main} guarantees this only for
    $q \gtrsim 10^4$--$10^5$, and the measured trajectory ($14$ at $q=211$, $19$ at $q=601$) is
    consistent with that timetable.
    The transitional effective exponent of $\xone$ over $101 \le q \le 601$, approximately
    $-0.76$, is the shadow of $\lvert B_{\none(q)}\rvert$ still growing while $q^{-2}$ pulls the
    ratio down; the asymptotic exponent is exactly $-2$.

    \begin{remark}[Prospective finite-$q$ resonance heuristic]
      \label{rem:prediction}
      At $q \in \{809, 1009, 1201\}$ the expected number of active resonance conditions per block
      near depth $22$ is still of order one.
      The resulting heuristic --- not a consequence of Theorem~\ref{thm:main} --- predicts
      $\none(q)<22$ there, with stabilisation at $\none=22$ only for
      $q \sim 10^4$--$10^5$.
      It is a pre-stated discriminating test for any future measurement campaign at those primes.
    \end{remark}

  \section{Degeneracy of the coverage-collapse ansatz}
    \label{sec:collapse}

    A natural structural route to the law of $\xone$ would have been the scaling collapse
    \begin{equation}
      r_q(n) \approx q\,\Phi\!\left(\frac{\Bball}{q^2}\right),
      \qquad \Phi \text{ smooth},\; \Phi(0)=0,\; \Phi \text{ increasing},
      \tag{$\star$}
      \label{eq:star}
    \end{equation}
    which would put the whole $q$-dependence of the rank trajectory into the reduced coverage
    $x = \Bball/q^2$ and reduce the threshold condition to a front equation
    $\Phi'(\xone) = \epssat\,q$.

    \begin{proposition}[The collapse is degenerate at the threshold front]
      \label{prop:collapse}
      By Theorem~\ref{thm:interval}, $r_c(n) \le 9(2n+3)$.
      Set $n=\lfloor\alpha\sqrt q\rfloor$ with fixed, sufficiently small $\alpha>0$ so that the
      finite ball is in the injective regime of Lemma~\ref{lem:isometry}.
      Then $r_c(n)=O(\sqrt q)=o(q)$ while
      $\lvert B_n\rvert/q^2\to\Cheis\alpha^4>0$.
      Hence~\eqref{eq:star} forces $\Phi$ to vanish on a non-trivial interval adjacent to zero.
      It cannot control the threshold front $\xone(q)\to0$: the natural support scale of the rank
      is $n/q$, not $\lvert B_n\rvert/q^2$.
    \end{proposition}

    The coverage variable $x$ retains its role as the definition of the observable
    $\xone$~\eqref{eq:x1}, but the law of $\xone$ is combinatorial (interval growth against exact
    sphere counts), not diffusive: as noted in Remark~\ref{prop:gap}, the spectral gap
    $\lambda_2$ never enters the rank dynamics.

  \section{Status and consequences}
    \label{sec:status}

    \begin{itemize}[nosep]
      \item \emph{Proved, deterministic:} the character reduction (Proposition~\ref{prop:character});
        the exact interval formula~\eqref{eq:interval}; the universal bounds
        $2 \le \Delta r_c(n) \le 2\lvert\Tp\rvert \le 18$; the shellwise resonance
        law~\eqref{eq:shellwise}; the strict crossing $n_*(10^{-3}) = 22$ on exact spheres;
        $\none(q) \le 22$ and $\xone(q) \asymp q^{-2}$ for $q \ge 311$
        (Theorem~\ref{thm:main}(i)); the window-closure Corollary~\ref{cor:window}.
      \item \emph{Proved, probabilistic:} $\none(q) \to 22$ and $q^2\xone(q) \to 99{,}689$ in
        probability under uniform generic block sampling (Theorem~\ref{thm:main}(ii)), with a
        conservative explicit bound (Remark~\ref{rem:conservative}).
      \item \emph{Empirical:} the measured transitional values of Table~\ref{tab:x1}, explained by
        the resonance mechanism of Theorem~\ref{thm:shellwise}.
      \item \emph{Not claimed:} almost-sure statements for the literal fixed-seed software
        sequence; sharp finite-$q$ probability constants; central-coordinate sensitivity of this
        observable; any mass-hierarchy mechanism.
    \end{itemize}

    \paragraph{Consequence for the spatial-limit programme.}
      In the dual-window limit of the canonical Fourier filtration~\cite{Beau2026q5a}, the
      potential coefficient is $\lambda = \lim_{q\to\infty} \xone(q)/\Cheis$.
      Theorem~\ref{thm:main}(i) gives $\lambda = 0$: the dual-window limit carries no oscillator
      potential and reduces to the free Dirichlet form.
      This closes the only identified conditional route from that harmonic confinement toward a
      spatial limit operator.
      It neither proves nor disproves hypothesis [H-L] of the emergent-geometry sub-programme by
      any other construction; the present result is spectral only.

    \subsection{Interpretive outlook}
      \label{sec:outlook}

      The following is a structural reading, not a result.
      The theorem says that the finite Weil--BFS observable has a \emph{finite novelty horizon}:
      beyond depth $22$ --- a constant fixed by the threshold and by exact sphere counts of the
      infinite Heisenberg group, independent of the arithmetic size $q$ of the substrate ---
      exploration creates no above-threshold density of new spectral directions per visited vertex.
      Novelty of this observable is a local, combinatorial phenomenon (growth of a sumset of at
      most nine progression classes), not a diffusive or spectral one; enlarging the universe of
      configurations deepens the transitional regime but never the horizon itself.
      Within the Cosmochrony programme this sharpens the reading of window-based capacity
      exponents: they are properties of a bounded relational neighbourhood of the origin, and
      their implemented extraction window closes once its proxy lower edge passes the horizon
      (Corollary~\ref{cor:window}).

    \paragraph{Data and Code Availability.}
      The window-depth measurement (\texttt{n1\_resumable.py}, multicore
      \texttt{n1\_parallel.py}), the ball-growth constant (\texttt{cheis\_ballgrowth.py},
      $\Cheis \approx 0.427$, $D = 4$), and the spectral-gap scaling
      (\texttt{lambda2\_scaling.py}) are in the Cosmochrony repository, directory
      \texttt{admissibility/o25/code}.
      Both depth scripts reproduce the O28 values $4, 8, 11, 13, 14$ at
      $q = 29, 61, 101, 151, 211$ exactly.
      The exact sphere counts of Table~\ref{tab:spheres} and the central-extent identity
      $\lfloor n^2/4\rfloor$ are reproduced by \texttt{code/h3z\_spheres.py} in this repository
      (exact integer BFS of $H_3(\mathbb{Z})$ to depth $31$).

    \paragraph{Acknowledgements.}
      AI-assisted tools were used during the preparation of this work.
      All scientific content, results, and conclusions are the sole responsibility of the author.

  \newpage
  \bibliographystyle{plain}
  \bibliography{cosmochrony-bibliography,external-refs}

\end{document}
